\footnotetext[1]{Corresponding author.}


\begin{abstract}
In this work, we focus on text-based person retrieval, which identifies individuals based on textual descriptions. Despite advancements enabled by synthetic data for pretraining, a significant domain gap—due to variations in lighting, color, and viewpoint—limits the effectiveness of the pretrain-finetune paradigm. To overcome this issue, we propose a unified pipeline incorporating domain adaptation at both image and region levels. Our method features two key components: Domain-aware Diffusion (DaD) for image-level adaptation, which aligns image distributions between synthetic and real-world domains, \eg, CUHK-PEDES, and Multi-granularity Relation Alignment (MRA) for region-level adaptation, which aligns visual regions with descriptive sentences, thereby addressing disparities at a finer granularity. This dual-level strategy effectively bridges the domain gap, achieving state-of-the-art performance on CUHK-PEDES, ICFG-PEDES, and RSTPReid datasets. 
The dataset, model, and code are available at \url{https://github.com/Shuyu-XJTU/MRA}.
\end{abstract}


% In this work, we focus on text-based person retrieval, which identifies individuals based on textual descriptions. Despite advancements enabled by synthetic data for pretraining, a significant domain gap—due to variations in lighting, color, and viewpoint—limits the effectiveness of the pretrain-finetune paradigm. To overcome this issue, we propose a unified pipeline incorporating domain adaptation at both image and region levels. Our method features two key components: Domain-aware Diffusion (DaD) for image-level adaptation, which aligns image distributions between synthetic and real-world domains, e.g., CUHK-PEDES, and Multi-granularity Relation Alignment (MRA) for region-level adaptation, which aligns visual regions with descriptive sentences, thereby addressing disparities at a finer granularity. This dual-level strategy effectively bridges the domain gap, achieving state-of-the-art performance on CUHK-PEDES, ICFG-PEDES, and RSTPReid datasets. The dataset, model, and code are available at https://github.com/Shuyu-XJTU/MRA.


\begin{figure}[!t]
\centering
\includegraphics[width=1\linewidth]{images/dataset.pdf}
\caption{Selected images from synthetic data generated by Diffusion model~\cite{yang2023towards}, real-world data, \ie, CUHK-PEDES~\cite{li2017person}, and our proposed Synthetic Domain-Aligned dataset (SDA).
We could observe that the visual gap between synthetic and real-world data (target domain) remains at illumination, color, viewpoints, \etc.
In contrast, images from SDA exhibit a target style while maintaining the high fidelity of the source image, characterized by a wide variety of variations in pose, appearance, background, \etc.
(Best viewed when zooming in.)
}
\label{fig: examples}
\end{figure}


